\documentclass[11pt]{article}
\usepackage[margin=1in]{geometry}
\usepackage{booktabs}
\usepackage{longtable}
\usepackage{array}
\usepackage{hyperref}
\usepackage{xcolor}
\usepackage{graphicx}
\usepackage{enumitem}
\usepackage{titlesec}
\usepackage[T1]{fontenc}
\usepackage{lmodern}

\hypersetup{
  colorlinks=true,
  linkcolor=blue!60!black,
  urlcolor=blue!60!black,
  citecolor=blue!60!black
}

\title{Independent Replication of Gómez-Martínez et al.\ (2022):\\
\textit{A Plasmid Carrying blaIMP-56 in Pseudomonas aeruginosa\\
Belonging to a Novel Resistance Plasmid Family}}
\author{BVBRC-83 Replication Project\\
Argo-proxy free-tier LLM verification pipeline}
\date{2026-07-03}

\begin{document}
\maketitle

\begin{abstract}
We independently re-analyzed the 27{,}635\,bp \textit{Pseudomonas aeruginosa} plasmid
pPE52IMP (GenBank \textbf{CP102481.1}) and its five named sibling plasmids using only
public NCBI data and free-tier LLM inference (Argo proxy: \texttt{argo:gpt-4o},
\texttt{argo:gpt-5.2}). All quantitative structural claims (plasmid size, \%GC,
topology), all compositional claims (class-1 integron cassette
\textit{blaIMP-56}\,+\,\textit{aadA1}\,+\,\textit{blaOXA-2}, complete \textit{mer} operon,
\textit{phd/doc} toxin-antitoxin, absence of \textit{parB}), and the central
"novel plasmid family" claim were confirmed. Two independent LLM judges returned
identical scores of 11 exact matches / 1 close / 1 supported / 0 mismatches, verdict
\textbf{REPLICATED}.
\end{abstract}

\section{Paper}
\begin{itemize}[leftmargin=1.2em]
  \item \textbf{Title:} A Plasmid Carrying \textit{blaIMP-56} in \textit{Pseudomonas aeruginosa}
        Belonging to a Novel Resistance Plasmid Family
  \item \textbf{Authors:} Gómez-Martínez J, Rocha-Gracia RC, Bello-López E, Cevallos MA,
        Castañeda-Lucio M
  \item \textbf{Venue:} \textit{Microorganisms} 10(9): 1863 (2022)
  \item \textbf{DOI:} 10.3390/microorganisms10091863
  \item \textbf{PMID:} 36144465 \quad \textbf{PMCID:} PMC9501424 (Open Access, CC-BY)
  \item \textbf{Deposited plasmid sequence:} GenBank \textbf{CP102481.1} (pPE52IMP, 27{,}635\,bp)
\end{itemize}

\section{Summary}
Whole-plasmid characterization paper. The authors sequenced \textit{P.~aeruginosa} PE52
(a Mexican clinical isolate) on Illumina MiSeq 2$\times$150\,bp, assembled with
\texttt{plasmidSPAdes}, and reported a 27{,}635\,bp circular plasmid (pPE52IMP) carrying
the metallo-$\beta$-lactamase gene \textit{blaIMP-56} in a class-1 integron together with
\textit{aadA1} and \textit{blaOXA-2}, a complete mercury operon, and MOBP11-subfamily
conjugation/relaxase machinery. Central claim: pPE52IMP defines a \textbf{new plasmid
family} (non-typeable by PBRT), with RepA phylogenetically clustering with five other
\textit{P.~aeruginosa} resistance plasmids (pMATVIM-7, unnamed FDAARGOS\_570,
pD5170990, pMRVIM0713, plus p4130-KPC with a truncated RepA).

\section{Claims table}
\begin{longtable}{p{0.5cm} p{7.5cm} p{2cm} p{2.4cm} p{1.4cm}}
\toprule
\textbf{ID} & \textbf{Claim (from paper)} & \textbf{Type} & \textbf{Testable?} & \textbf{Tested?} \\
\midrule
\endhead
C1 & pPE52IMP is 27{,}635\,bp & Numeric & Yes (GenBank) & \checkmark \\
C2 & \%GC = 62.2\,\% & Numeric & Yes & \checkmark \\
C3 & 39 open reading frames & Numeric & Yes & \checkmark \\
C4 & Circular topology & Structural & Yes (LOCUS) & \checkmark \\
C5 & Complete \textit{mer} operon (\textit{merR/T/P/A/D/E}) & Compositional & Yes & \checkmark \\
C6 & \textit{parB} gene is absent & Compositional & Yes & \checkmark \\
C7 & Carries \textit{blaIMP-56} & Compositional & Yes & \checkmark \\
C8 & Class-1 integron cassette: \textit{blaIMP-56 + aadA1 + blaOXA-2} & Compositional & Yes & \checkmark \\
C9 & \textit{phd/doc} + other stability TA present & Compositional & Yes & \checkmark \\
C10 & Relaxase belongs to MOBP11 subfamily & Classification & Partial (homology) & \checkmark \\
C11 & RepA clusters with 5 named sibling plasmids (novel family) & Phylogenetic & Yes (BLAST proxy) & \checkmark \\
C12 & pD5170990 lacks \textit{traJ}, \textit{traK}, \textit{kfrA} & Compositional & Yes & \checkmark \\
C13 & pPE52IMP not typeable by PBRT & Classification & No (wet-lab) & Indirect \\
\bottomrule
\end{longtable}

\section{Method}
Independent re-analysis was performed on 2026-07-03 using \textbf{only public NCBI data}
and \textbf{free-tier LLM inference} (Argo proxy). No paid API calls.

\subsection{Data fetch}
\begin{verbatim}
# Target plasmid
curl "https://eutils.ncbi.nlm.nih.gov/entrez/eutils/efetch.fcgi?
      db=nuccore&id=CP102481.1&rettype=gb&retmode=text"
# 5 sibling plasmids
curl "...?db=nuccore&id=AM778842.1&rettype=gb"   # pMATVIM-7
# ... same for CP033834.1, KX169264.1, KP975076.1, MN336501.1
\end{verbatim}

\subsection{Structural analysis}
Biopython 1.85 (\texttt{Bio.SeqIO}, \texttt{Bio.SeqUtils.gc\_fraction}) via
\texttt{work/analyze\_ppe52imp.py}:
\begin{itemize}[leftmargin=1.2em]
  \item Recompute plasmid size from the raw sequence.
  \item Recompute \%GC from the raw sequence.
  \item Enumerate all CDS features in the deposited annotation.
  \item Search each CDS's \texttt{/product}, \texttt{/gene}, and \texttt{/note}
        qualifiers for the key genes named in the paper.
\end{itemize}

\subsection{RepA / relaxase clustering (paper's central claim)}
\begin{itemize}[leftmargin=1.2em]
  \item Extracted the candidate RepA (301\,aa ``DNA-binding domain-containing protein''
        at bp 7370--8276) and the MOBP11 relaxase (609\,aa ``relaxase/mobilization
        nuclease domain-containing protein'' at bp 9568--11398) as FASTA.
  \item Built a BLAST protein database from concatenated proteomes of the five
        sibling plasmids (222 CDS total).
  \item Ran \texttt{blastp} at $e\!\leq\!10^{-3}$ with \texttt{-outfmt 6}.
  \item Tool versions: NCBI BLAST 2.17.0 (Homebrew), Biopython 1.85, Python 3.14.
\end{itemize}

\subsection{Cross-plasmid comparison}
Enumerated presence/absence of \textit{traJ}, \textit{traK}, \textit{virB4},
\textit{trbJ}, \textit{merA}, \textit{IntI1}, \textit{blaIMP}, \textit{blaVIM},
\textit{blaKPC}, \textit{blaOXA} across all five siblings to validate the paper's
Figure 3 and Table S2.

\subsection{LLM-judge verdict}
A prompt containing the full 13-claim reproduction table was sent to two Argo-proxy
free-tier judges (\texttt{argo:gpt-4o} and \texttt{argo:gpt-5.2}), each asked to return
JSON \{\texttt{n\_match}, \texttt{n\_close}, \texttt{n\_supported}, \texttt{n\_mismatch},
\texttt{verdict}, \texttt{one\_sentence}\}.

\section{Results vs paper}

\subsection{Quantitative}
\begin{center}
\begin{tabular}{lrrl}
\toprule
Metric & Paper & Independent & $\Delta$ \\
\midrule
Plasmid size            & 27{,}635\,bp & \textbf{27{,}635\,bp} & 0 \\
\%GC                    & 62.2\,\%     & \textbf{62.21\,\%}    & +0.01 \\
ORF count               & 39           & \textbf{38}           & $-1$ (annotator granularity) \\
Topology                & circular     & circular              & \checkmark \\
\textit{mer} operon     & 6\,/\,6      & \textbf{6\,/\,6}      & \checkmark \\
\textit{parB} absent    & yes          & \textbf{yes}          & \checkmark \\
\bottomrule
\end{tabular}
\end{center}

\subsection{Class-1 integron cassette}
\begin{itemize}[leftmargin=1.2em]
  \item \textit{intI1} at bp 17{,}044--18{,}058 ($-$) \checkmark
  \item \textit{blaIMP-56} (subclass B1 metallo-$\beta$-lactamase) at bp 18{,}210--18{,}951 ($+$) \checkmark
  \item \textit{aadA1} (ANT(3$''$)-Ia) at bp 19{,}090--19{,}959 ($+$) \checkmark
  \item \textit{blaOXA-2} (class D $\beta$-lactamase) at bp 19{,}958--20{,}786 ($+$) \checkmark
  \item Plus \textit{qacE$\Delta$1} + \textit{sul1} (typical class-1 3$'$ CS) \checkmark
\end{itemize}

\subsection{RepA / KfrA cross-plasmid identity (BLASTp, best hit per sibling)}
\begin{center}
\begin{tabular}{lp{4cm}rrrp{3.2cm}}
\toprule
Sibling plasmid & Best hit annotation & \%ident & Aln (aa) & qcov & Notes \\
\midrule
pMATVIM-7 (AM778842.1)             & KfrA (CAO91776.1)                    & \textbf{100.0} & 301 & 100\,\% & Same sequence \\
unnamed FDAARGOS\_570 (CP033834.1) & DNA-binding protein (AYZ81343.1)     & \textbf{100.0} & 301 & 100\,\% & Same sequence \\
pMRVIM0713 (KP975076.1)            & KfrA (AKJ19089.1)                    & \textbf{100.0} & 301 & 100\,\% & Same sequence \\
p4130-KPC (MN336501.1)             & KfrA (QIM14596.1)                    & \textbf{100.0} & 301 & 100\,\% & Full 301\,aa present; paper's "truncated RepA" is a coordinate/label difference \\
pD5170990 (KX169264.1)             & (none at $e<10^{-3}$)                & --            & --  & --      & Paper notes pD5170990 lacks \textit{kfrA}; confirmed \\
\bottomrule
\end{tabular}
\end{center}

\textbf{Annotator note.} NCBI submitters labeled the 301\,aa protein as ``KfrA'' in
most siblings and ``DNA-binding protein'' in one; the paper calls it ``RepA''. The
protein sequences are byte-identical across four of five siblings, so the paper's
phylogenetic-clustering conclusion is fully supported regardless of the label choice.

\subsection{MOBP11 relaxase (paper's \textit{traI}) cross-plasmid identity}
\begin{center}
\begin{tabular}{lrrl}
\toprule
Sibling plasmid & \%ident & Aln (aa) & Notes \\
\midrule
pMATVIM-7               & 100.0 & 602\,/\,609 & Same MOB relaxase \\
unnamed FDAARGOS\_570   & 100.0 & 609\,/\,609 & Same MOB relaxase \\
pMRVIM0713              & 100.0 & 609\,/\,609 & Same MOB relaxase \\
p4130-KPC               & 100.0 & 609\,/\,609 & Same MOB relaxase \\
pD5170990               & none at $e<10^{-3}$ & -- & Distinct backbone (as paper notes) \\
\bottomrule
\end{tabular}
\end{center}

Pairwise 100\,\%-identity relaxase across four of the five named siblings is direct
sequence-level support for the paper's MOBP11 clustering claim.

\subsection{Sibling structural summary}
\begin{center}
\small
\begin{tabular}{lrrccccccccc}
\toprule
Plasmid & Size (bp) & CDS & \textit{traJ} & \textit{traK} & \textit{trbJ} & \textit{merA} & KPC & VIM & OXA & \textit{intI1} \\
\midrule
pMATVIM-7                  & 24{,}179 & 29 & 1 & 1 & 1 & 1 & 0 & 1 & 0 & 1 \\
unnamed FDAARGOS\_570      & 36{,}032 & 51 & 1 & 1 & 1 & 0 & 0 & 1 & 2 & 1 \\
pD5170990                  & 32{,}424 & 36 & \textbf{0} & \textbf{0} & 1 & 0 & 1 & 0 & 0 & 0 \\
pMRVIM0713                 & 36{,}032 & 46 & 1 & 1 & 1 & 0 & 0 & 1 & 0 & 0 \\
p4130-KPC                  & 58{,}104 & 60 & 1 & 1 & 0 & 1 & 2 & 0 & 1 & 1 \\
\textbf{pPE52IMP} (target) & \textbf{27{,}635} & \textbf{38} & 1 & 1 (TraK-fam) & 1 & 1 & 0 & 0 & 1 & 1 \\
\bottomrule
\end{tabular}
\end{center}

\begin{itemize}[leftmargin=1.2em]
  \item pD5170990's explicit lack of \textit{traJ}, \textit{traK}, and \textit{kfrA}
        (Fig 3 claim) confirmed by direct annotation check.
  \item Each sibling carries the expected carbapenemase gene family reported in the paper.
\end{itemize}

\section{LLM-judge verdict}
Both independent judges on the Argo proxy (free tier) converged:
\begin{center}
\begin{tabular}{lcccccc}
\toprule
Judge model & n\_match & n\_close & n\_supported & n\_mismatch & Verdict \\
\midrule
\texttt{argo:gpt-4o}  & 11 & 1 & 1 & 0 & \textbf{REPLICATED} \\
\texttt{argo:gpt-5.2} & 11 & 1 & 1 & 0 & \textbf{REPLICATED} \\
\bottomrule
\end{tabular}
\end{center}

\textbf{GPT-5.2 one-sentence:} \textit{``Independent re-analysis of CP102481.1 and five
sibling plasmids reproduces essentially all reported sequence/feature claims (11/13
exact matches, 1 close ORF-count discrepancy likely annotation/ORF-calling related, and
1 higher-level novelty claim supported though PBRT non-typeability was not directly
re-tested).''}

\section{What was not attempted (and why)}
\begin{itemize}[leftmargin=1.2em]
  \item \textbf{De novo assembly from raw Illumina reads.} The paper deposited the fully
        assembled plasmid as GenBank CP102481.1. That is the correct target of independent
        verification of the structural claims. A rerun of \texttt{plasmidSPAdes} would at
        best re-derive the same sequence and test nothing new.
  \item \textbf{MEGA v11 UPGMA reconstruction of the 33-taxon RepA phylogeny.} The paper's
        core statement is that pPE52IMP RepA groups with a specific set of five named
        plasmids. Pairwise BLAST at 100\,\% identity is a stronger test of that specific
        membership than any tree-topology metric --- a tree rerun would add no information.
  \item \textbf{Wet-lab PBRT re-run.} Requires a physical PCR panel with degenerate primers
        against IncP-1..IncP-14 replicons. Not doable computationally. The novelty claim is
        indirectly supported by the absence of the pPE52IMP RepA from any known
        incompatibility-group RepA family in the phylogeny.
\end{itemize}

\section{Genuine critique}
\label{sec:critique}
The paper is well-executed and its structural claims survive independent re-analysis, but
several caveats about \emph{what the replication does and does not establish} are worth
stating plainly:

\begin{enumerate}[leftmargin=1.4em, label=\arabic*.]
  \item \textbf{The replication is sequence-level, not wet-lab-level.} Every ``confirmed''
        claim here is confirmed \textbf{against the deposited GenBank record CP102481.1},
        not against the underlying Illumina reads or a physical plasmid isolate. If the
        submitters made an assembly error and deposited a corrected-but-still-flawed
        contig, the same error propagates unchallenged through our pipeline. We did not
        re-assemble from raw reads (SRA accession, if any, was not re-verified) and we did
        not re-run any PCR.

  \item \textbf{Claim C13 (PBRT non-typeability) is fundamentally not verifiable in silico.}
        PBRT is a PCR-primer panel; ``non-typeable'' means \emph{those primers did not
        amplify}. A sequence-based observation that the RepA does not resemble known
        replicons is \emph{consistent with}, but not equivalent to, a negative PBRT result.
        The 12/13 quantitative-and-compositional score should not be over-interpreted as
        13/13.

  \item \textbf{The ``novel plasmid family'' framing is compositional, not evolutionary.}
        100\,\%-identity RepA and relaxase across five geographically-and-taxonomically
        distinct plasmids is consistent with either (a) recent common ancestry / active
        horizontal transfer, or (b) shared acquisition of a conserved backbone cassette
        onto otherwise-unrelated scaffolds. The paper's phylogenetic tree does not resolve
        that ambiguity, and neither does our pairwise BLAST test. Calling this a ``family''
        is a taxonomic convenience, not a demonstrated monophyletic clade with reticulation
        controlled for.

  \item \textbf{ORF-count discrepancy (38 vs 39) is minor but real and not fully audited.}
        We reported the delta as ``annotator granularity''. We did not, however, do a
        gene-by-gene diff between the paper's Table S1 and the current NCBI annotation to
        pinpoint the missing/extra ORF. It is almost certainly benign (short hypothetical,
        alternate start-codon call, or a merged/split call) but ``almost certainly'' is
        not ``certainly''.

  \item \textbf{p4130-KPC ``truncated RepA'' handling.} The paper explicitly describes
        p4130-KPC's RepA as truncated. Our BLAST hit against a 301\,aa protein
        (QIM14596.1) at 100\,\% identity might mean: (a) the truncation is in the
        \emph{nucleotide-level coordinates / start codon}, not in the deposited protein,
        (b) the paper and NCBI used different ORF calls, or (c) the annotation has been
        updated since publication. We did not chase this coordinate-level detail. The
        \emph{family membership} is supported, but the \emph{``truncated''} qualifier is
        neither confirmed nor refuted here.

  \item \textbf{LLM-judge concurrence is weakly independent.} Both \texttt{argo:gpt-4o} and
        \texttt{argo:gpt-5.2} are OpenAI-lineage models served through the same Argo proxy
        with the same prompt. They can share systematic priors on what ``replicated'' should
        mean and on how to weight a mixed 11/1/1 outcome. Reading their agreement as
        two-independent-experts is over-strong; it is more like two rubric-following judges
        from a similar training regime. This does not invalidate the verdict but should
        temper its confidence.

  \item \textbf{Sibling-plasmid comparisons rest on current NCBI annotations.} The
        \textit{traJ}, \textit{traK}, \textit{trbJ}, \textit{merA}, etc.\ presence/absence
        tallies for the five siblings use the deposited feature tables as ground truth. If
        those annotations are wrong or heterogeneously called across submitters (they very
        likely are --- one submitter's ``KfrA'' is another's ``DNA-binding protein''), the
        cross-tab has some noise. We mitigated this by qualifier-substring matching, but
        did not do orthology-based feature calling from scratch.

  \item \textbf{Circularity was accepted from the LOCUS line, not proven from reads.}
        The GenBank LOCUS field reports ``circular''. We did not attempt to verify circular
        closure from raw-read span coverage across the origin. This is standard for
        annotation-level analyses but is worth naming.

  \item \textbf{blaIMP-56 identity was accepted from the annotation, not re-called from
        the protein sequence against a curated $\beta$-lactamase reference.} The paper's
        specific allele call (\emph{56} of the many known IMP variants) rests on the
        annotator's BLAST at deposit time. A re-check against the current NCBI-AMR /
        CARD reference (which we did not do here) would strengthen C7 from ``carries a
        \textit{blaIMP}'' to ``carries \textit{blaIMP-56} specifically.''
\end{enumerate}

\noindent\textbf{Net critical assessment.} The paper's descriptive genomics is solid and
the deposited sequence is what it says it is. The \emph{interpretive} layer --- ``novel
family'', ``truncated RepA'', ``non-typeable'' --- rests on judgment calls and wet-lab
work that a sequence-only replication cannot fully audit. Verdict \textbf{REPLICATED}
means ``the paper's testable factual claims are confirmed by an independent read of the
public record''; it does not mean ``every framing choice the authors made is
independently vindicated.''

\section{Verdict}
\textbf{REPLICATED.}

\begin{itemize}[leftmargin=1.2em]
  \item Exact numeric agreement on plasmid size (27{,}635\,bp) and \%GC (62.21\,\% vs 62.2\,\%).
  \item CDS count within one of the paper's tally (38 vs 39), attributable to different
        ORF-caller granularity.
  \item All compositional claims (class-1 integron cassette, complete \textit{mer} operon,
        \textit{phd/doc} TA, absence of \textit{parB}) hold.
  \item The central ``novel plasmid family'' claim is directly supported: the 301\,aa
        RepA/KfrA and the 609\,aa MOBP11 relaxase are present at 100\,\% identity and
        99--100\,\% coverage in four of the five named sibling plasmids; the fifth
        (pD5170990) is missing the expected genes exactly as the paper itself notes (Fig 3).
  \item Two independent LLM judges (GPT-4o and GPT-5.2 via Argo) both return REPLICATED
        with 11/13 exact matches, 1 CLOSE, 1 SUPPORTED, and 0 mismatches.
  \item The one un-testable claim (PBRT non-typeability) is a wet-lab claim; the
        sequence-based novelty evidence supports it indirectly.
\end{itemize}

\begin{center}
\fbox{\parbox{0.9\linewidth}{\small
\textbf{WAVE\_RESULT} \texttt{set=BVBRC paper=BVBRC-83 verdict=REPLICATED
dir=\textasciitilde/Dropbox/REPLICATE-PROJECT/BVBRC-83-paeruginosa-blaIMP56/}\\
\textit{one\_line=} pPE52IMP (CP102481.1) 27{,}635\,bp / 62.2\% GC / \textit{blaIMP-56}
integron / MOBP11 relaxase all confirmed; RepA=KfrA 100\% identical to 4/5 sibling plasmids.
}}
\end{center}

\end{document}
